\section*{Chapitre \ref{ch:nw:api-engineering-for-crypto-venues} --- Ingénierie des API pour les plateformes crypto}

\begin{solution}{exo:nw:api-engineering-for-crypto-venues:1}
$6\,000/250 = 24$, et rien d'autre pendant cette minute.
\end{solution}

\begin{solution}{exo:nw:api-engineering-for-crypto-venues:2}
Pour limiter ce que coûtent les problèmes d'une connexion~: le \omterm{def:nw:networking-for-trading:hol}{blocage en tête de file} derrière un symbole très actif, les carnets perdus quand elle tombe, les instantanés nécessaires pour
se rétablir, et la charge de la réinitialisation de 24 heures. Cinq connexions de 160 se partagent le risque et le travail.
\end{solution}

\begin{solution}{exo:nw:api-engineering-for-crypto-venues:3}
Répondre par un pong reprenant la charge utile du ping, dès que possible (RFC 6455). La plateforme envoie un ping toutes les 20 secondes et déconnecte une connexion qui
n'a pas répondu dans la minute.
\end{solution}

\begin{solution}{exo:nw:api-engineering-for-crypto-venues:4}
En 10 secondes, une adresse peut ouvrir $300 \times 10/300 = 10$ connexions selon la lecture d'une répartition uniforme, donc 40 connexions exigent 4 adresses.
\end{solution}

\begin{solution}{exo:nw:api-engineering-for-crypto-venues:5}
La réception moins l'événement est le délai dans un sens moins le \omterm{def:nw:time-synchronisation:offset}{décalage d'horloge} de la plateforme~; le délai a un plancher que certains messages atteignent dans chaque fenêtre, donc le
minimum est le plancher moins le décalage, et ses variations sont celles du décalage. Il échoue lorsque le plancher du délai se déplace (changement de route, congestion qui dure toute la
fenêtre) ou lorsque la plateforme horodate les événements à des points différents selon les flux.
\end{solution}

\begin{solution}{exo:nw:api-engineering-for-crypto-venues:6}
Cette mesure mélange deux horloges~: l'écart inclut le \omterm{def:nw:time-synchronisation:offset}{décalage d'horloge} de la plateforme (en millisecondes dans la simulation du chapitre) et le délai à sens unique vers la firme.
Mesurer le tick-to-order sur l'horloge propre de la firme à partir de l'horodatage d'arrivée, et n'utiliser l'heure de la plateforme qu'après avoir retranché la dérive d'horloge estimée.
\end{solution}

\begin{solution}{exo:nw:api-engineering-for-crypto-venues:7}
75~: $80 \times 75 = 6\,000$, tout le budget de la minute~; à 76, le dernier instantané attend la minute suivante.
\end{solution}

\begin{solution}{exo:nw:api-engineering-for-crypto-venues:8}
Une connexion est un domaine de défaillance~: une déconnexion ou la réinitialisation de 24 heures fait tomber tous les carnets, un symbole très actif retarde tous les autres, et la reprise par
connexion exige un instantané pour chaque symbole, ce que le budget de poids ne peut pas absorber assez vite. Davantage de connexions, chacune sous le cap, sont plus fiables.
\end{solution}

\begin{solution}{pb:nw:api-engineering-for-crypto-venues:1}
\textbf{Partie I.}
\begin{enumerate}
\item 800 flux, 5 connexions de 160, 1 adresse.
\item Une connexion.
\item Toujours une~: cinq reconnexions en 10 secondes restent largement sous la limite par adresse.
\item Chaque connexion doit être remplacée au moins une fois par jour~: il faut échelonner les réinitialisations et se reconnecter avant la plateforme, une nouvelle connexion étant abonnée avant
  que l'ancienne ne se ferme.
\end{enumerate}
\textbf{Partie II.}
\begin{enumerate}[resume]
\item Cinq (les messages perdus touchent cinq symboles) par symbole~; les 80 de la connexion, par connexion.
\item \qty{180.1}{\second}, et 5\,769 secondes-symbole périmées.
\item \qty{0.2}{\second} et 11.2 secondes-symbole~: 80 instantanés de poids 50 tiennent dans une minute.
\item Envoyer toute requête qui demande du poids, ordres compris, jusqu'à ce que les instantanés de la reprise aient épuisé leurs fenêtres.
\end{enumerate}
\textbf{Partie III.}
\begin{enumerate}[resume]
\item Risk 500 (tout ce qu'il demande), market making 3\,300, arbitrage 2\,200.
\item Comme une part réservée, ou via la classe risque~: la reprise ne doit pas être privée de ressources par le trading, ni en priver celui-ci.
\item À moins de \qty{15}{\micro\second} de la valeur vraie, de \qty{-2.23}{\milli\second} au début à \qty{-4.20}{\milli\second} à la fin.
\item Mesurer la latence dans les deux sens avec la plateforme, sur une même base de temps, et repérer quand l'horloge de la plateforme, donc ses horodatages, se déplace.
\end{enumerate}
\textbf{Partie IV.}
\begin{enumerate}[resume]
\item \emph{Résultat nommé}~: 800 flux exigent cinq connexions de 160 et une adresse sous les limites de la plateforme~; cinq messages perdus coûtent 0.5
  secondes-symbole périmées avec la reprise par symbole, et 5\,769 secondes-symbole (le dernier carnet revenu après trois minutes) avec la reprise par connexion, à partir de
  d'instantanés de pleine profondeur.
\item Reprise par symbole~; puis la profondeur des instantanés.
\item La firme ne pouvait pas savoir quels symboles avaient manqué des messages et devrait reconstruire chaque carnet de la connexion~: des répartitions plus fines et des instantanés moins
  profonds seraient alors les seules parades.
\item De façon régulière~: quelques connexions à la fois, chacune renouvelée avant que la plateforme ne la ferme.
\item Une seconde copie de chaque événement~: un message manqué sur l'une est en général présent sur l'autre, donc le carnet n'a pas besoin d'instantané~; la première copie l'emporte.
\item Fenêtres fixes alignées sur celles de la plateforme, plusieurs règles à la fois (poids, ordres par intervalle, par jour), et la pause après les réponses 429 et 418.
\item Dans les lignes crypto~: répartitions, adresses, régulateurs par limite, conception de la reprise et profondeur des instantanés, et surveillance de la dérive d'horloge.
\item Soit quelques dixièmes de seconde pour un seul carnet, soit des minutes pour chaque carnet d'une connexion et pour tout le budget de débit de la firme.
\end{enumerate}
\end{solution}

\begin{solution}{iq:nw:api-engineering-for-crypto-venues:1}
Mettre les deltas en tampon, récupérer un instantané, écarter les deltas dont l'identifiant de mise à jour est antérieur ou égal à celui de l'instantané, appliquer le reste dans l'ordre~; si le premier identifiant de mise à jour d'un delta dépasse
le dernier identifiant du carnet plus un, abandonner le carnet et recommencer~; vérifier avec la somme de contrôle de la plateforme lorsqu'elle en propose une.
\iqlookfor{Numéros de séquence~; règle des trous~; ordre de l'instantané et du tampon~; sommes de contrôle.}
\end{solution}

\begin{solution}{iq:nw:api-engineering-for-crypto-venues:2}
Sous le cap de la plateforme par connexion (quelques centaines de flux chacune), une douzaine de connexions réparties sur deux adresses ou plus, des symboles équilibrés selon le
débit de messages, reconnexions échelonnées, reprise par symbole et une connexion de réserve prête.
\iqlookfor{Caps et limites~; domaines de défaillance~; échelonnement~; conception de la reprise.}
\end{solution}

\begin{solution}{iq:nw:api-engineering-for-crypto-venues:3}
Faire passer chaque requête par un régulateur par limite de la plateforme et attribuer à chaque stratégie une part selon sa priorité (water-filling), avec une part réservée au risque et
à la reprise~; refuser ou mettre en file ce qui dépasse une part plutôt que de laisser la plateforme répondre par une erreur 429.
\iqlookfor{Régulateur central~; allocation~; réserves~; aucune dépendance aux erreurs de la plateforme.}
\end{solution}

\begin{solution}{iq:nw:api-engineering-for-crypto-venues:4}
Horodater les arrivées avec une horloge synchronisée et suivre le minimum de l'écart entre l'arrivée et l'heure d'événement de la plateforme sur des fenêtres~: son niveau est le délai minimal moins le
\omterm{def:nw:time-synchronisation:offset}{décalage d'horloge} de la plateforme, et ses variations sont celles du décalage. Recouper avec les allers-retours mesurés grâce au point de terminaison de l'heure de la plateforme.
\iqlookfor{Filtre de délai minimal~; séparation du délai et du décalage~; vérification.}
\end{solution}

\begin{solution}{iq:nw:api-engineering-for-crypto-venues:5}
Selon la latence, la reprise et les limites~: une session WebSocket ou une session FIX persistante évite les poignées de main par requête et fournit des numéros de séquence (FIX) ou des
accusés de réception poussés~; l'interface REST est simple et sans état, mais plus lente~; il faut vérifier si la plateforme compte tous les chemins sous une même limite et laquelle propose des \omterm{def:nw:private-connectivity-to-crypto-venues:endpoint}{points de terminaison privés}.
\iqlookfor{Session contre requête~; sémantique de la reprise~; limites partagées.}
\end{solution}

\begin{solution}{iq:nw:api-engineering-for-crypto-venues:6}
La réinitialisation de 24 heures ou un événement planifié de la plateforme (connexions toutes ouvertes ensemble, réinitialisées ensemble)~; le calendrier propre à la machine (une tâche cron, le
ramasse-miettes, une sauvegarde)~; les taux de messages à ce moment-là (un rééquilibrage d'indice, une échéance de financement)~; et les limites (5 messages entrants par seconde, limites de
reconnexion) atteintes par une rafale de réabonnements.
\iqlookfor{Causes périodiques~; échelonnement~; limites~; mesurer avant de deviner.}
\end{solution}
